\documentclass[11pt,a4paper,fontset=fandol]{ctexart}
\usepackage[a4paper,top=2.3cm,bottom=2.5cm,left=2.3cm,right=2.3cm,headheight=18pt,headsep=0.55cm,footskip=25pt]{geometry}
\usepackage{amsmath,amssymb,amsthm,fancyhdr,lastpage,enumitem,titlesec,xcolor,hyperref,bookmark}
\usepackage[most]{tcolorbox}
\usepackage{setspace,array}
\hypersetup{colorlinks=true,linkcolor=blue!50!black,urlcolor=blue!50!black}
\setstretch{1.15}
\setlength{\parindent}{2em}
\setlength{\parskip}{0.28em}
\allowdisplaybreaks
\setlist[enumerate]{itemsep=0.42em,topsep=0.35em,leftmargin=2.35em}
\setlist[itemize]{itemsep=0.25em,topsep=0.25em,leftmargin=1.9em}
\definecolor{mainblue}{RGB}{36,83,140}
\definecolor{lightblue}{RGB}{241,246,252}
\definecolor{lightgray}{RGB}{246,246,246}
\definecolor{accent}{RGB}{153,76,0}
\definecolor{rulegray}{RGB}{110,118,128}
\titleformat{\section}{\Large\bfseries\color{mainblue}}{\thesection.}{0.45em}{}[\vspace{0.2em}\color{rulegray}\titlerule]
\titlespacing*{\section}{0pt}{1.15em}{0.75em}
\titleformat{\subsection}{\large\bfseries\color{mainblue}}{\thesubsection}{0.5em}{}
\newtcolorbox{infobox}[1]{breakable,colback=lightblue,colframe=mainblue,boxrule=0.7pt,arc=1mm,title=\textbf{#1},fonttitle=\bfseries}
\newtcolorbox{notebox}[1]{breakable,colback=lightgray,colframe=black!50,boxrule=0.55pt,arc=1mm,title=\textbf{#1},fonttitle=\bfseries}
\newtheoremstyle{problemstyle}{0.8em}{0.8em}{\normalfont}{}{\bfseries\color{accent}}{．}{0.6em}{}
\theoremstyle{problemstyle}
\newtheorem{exercise}{题目}[section]
\newenvironment{solution}{\par\noindent\textbf{解：}}{\hfill$\square$\par}
\newenvironment{testsolution}{\par\noindent\textbf{参考解答：}}{\hfill$\square$\par}
\newcommand{\answerspace}{\vspace{2.1cm}}
\newcommand{\biganswerspace}{\vspace{3cm}}
\newcommand{\fl}[1]{\left\lfloor #1\right\rfloor}
\newcommand{\fr}[1]{\left\{#1\right\}}

\pagestyle{fancy}\fancyhf{}
\fancyhead[L]{\small 取整与小数部分提高训练题单}
\fancyhead[C]{\small 代数专题：方程、不等式与证明}
\fancyhead[R]{\small 刘文博}
\fancyfoot[C]{\small 第 \thepage 页 / 共 \pageref{LastPage} 页}
\renewcommand{\headrulewidth}{0.55pt}\renewcommand{\footrulewidth}{0.35pt}

\begin{document}
\thispagestyle{empty}
\begin{titlepage}\centering\vspace*{0.9cm}
\begin{tcolorbox}[enhanced,width=0.92\textwidth,colback=white,colframe=mainblue,boxrule=1pt,arc=0mm,left=6mm,right=6mm,top=8mm,bottom=8mm]
\centering
{\fontsize{24}{30}\selectfont\bfseries 取整函数与小数部分提高训练题单\par}
\vspace{0.6cm}{\fontsize{17}{22}\selectfont 面向初中数学联赛的分层训练\par}
\vspace{1cm}
\begin{tcolorbox}[colback=lightblue,colframe=mainblue,width=0.84\textwidth,boxrule=1pt]
\centering
{\large 从基本计算推进到方程、不等式、恒等式与整数参数}\\[0.35em]
{\normalsize 建议分 3 次完成，先独立作答，再用详细解答订正}
\end{tcolorbox}
\vspace{1cm}
\renewcommand{\arraystretch}{1.42}
\begin{tabular}{>{\bfseries}p{3.4cm}p{8cm}}
题目数量： & 18 题，按难度递进 \\
核心方法： & 设 $x=n+t$、区间翻译、整数筛选、误差估计 \\
答案形式： & 末尾附详细解答和关键方法说明 \\
建议用时： & 每组 45--60 分钟
\end{tabular}
\vspace{0.9cm}{\large \today\par}
\end{tcolorbox}\end{titlepage}
\tableofcontents\newpage

\section{使用建议}
\begin{infobox}{每道题先做三个动作}
先写 $[u]\le u<[u]+1$；若同一个变量反复出现，再设 $x=n+t$，其中 $n\in\mathbb Z$、$0\le t<1$；最后单独检查整数条件和端点。
\end{infobox}
\begin{notebox}{训练节奏}
第 1--6 题巩固计算与化简；第 7--12 题集中训练方程、不等式；第 13--18 题进入联赛常见的证明、求和与参数问题。
\end{notebox}

\section{计算与结构识别}
\begin{exercise}计算 $[-5.02]+\{-5.02\}$，以及 $[\sqrt{17}]+[-\sqrt{17}]$。\end{exercise}
\begin{exercise}已知 $[x]=3$，求 $[4x]$ 的所有可能值，并说明每个值都能取得。\end{exercise}
\begin{exercise}求 $[x]+[-x]+\{x\}+\{-x\}$ 的所有可能值。\end{exercise}
\begin{exercise}证明并使用公式
\[
[2x]=2[x]+[2\{x\}]
\]
计算 $[2x]-2[x]$ 的所有可能值及对应条件。\end{exercise}
\begin{exercise}化简
\[
[x]+[x+\tfrac14]+[x+\tfrac24]+[x+\tfrac34].
\]\end{exercise}
\begin{exercise}设 $n$ 为正整数，求
\[
\sum_{k=0}^{n-1}\left[\frac{x+k}{n}\right]
\]
与 $[x]$ 的关系。\end{exercise}

\section{方程与不等式}
\begin{exercise}解方程 $2[x]-x=1$。\end{exercise}
\begin{exercise}解方程 $[x]+[x+\frac12]=5$。\end{exercise}
\begin{exercise}解方程 $[2x]=[x]+3$。\end{exercise}
\begin{exercise}解方程 $\{x\}=\frac{x+1}{3}$。\end{exercise}
\begin{exercise}解不等式 $1<[3x-2]\le6$。\end{exercise}
\begin{exercise}解不等式 $[x]<3\{x\}$。\end{exercise}

\section{证明与综合提高}
\begin{exercise}证明：对任意实数 $x,y$，
\[
[x]+[y]\le[x+y]\le[x]+[y]+1,
\]
并分别写出两个等号成立的条件。\end{exercise}
\begin{exercise}证明：对任意实数 $x$ 和正整数 $n$，
\[
[nx]=\sum_{k=0}^{n-1}\left[x+\frac{k}{n}\right].
\]\end{exercise}
\begin{exercise}证明：对任意实数 $x$ 和正整数 $n$，
\[
\left[\frac{[x]}{n}\right]=\left[\frac{x}{n}\right].
\]\end{exercise}
\begin{exercise}设 $x+y+z\in\mathbb Z$，证明 $\{x\}+\{y\}+\{z\}$ 只能取 $0,1,2$。\end{exercise}
\begin{exercise}求所有实数 $x$，使得
\[
[x]+[2x]+[4x]=7x.
\]\end{exercise}
\begin{exercise}求所有正整数 $n$，使得对一切实数 $x$ 都有
\[
[nx]=n[x].
\]\end{exercise}

\newpage
\section{详细解答}
\subsection{计算与结构识别}
\begin{solution}第 1 题：$[-5.02]=-6$，$\{-5.02\}=0.98$，和为 $-5.02$。又 $4<\sqrt{17}<5$，所以 $[\sqrt{17}]=4$、$[-\sqrt{17}]=-5$，和为 $-1$。\end{solution}
\begin{solution}第 2 题：设 $x=3+t$，$0\le t<1$，则 $[4x]=12+[4t]$，可能为 $12,13,14,15$。分别取 $t=0,\frac14,\frac12,\frac34$ 即可取得。\end{solution}
\begin{solution}第 3 题：若 $x\in\mathbb Z$，四项之和为 $0$；若 $x\notin\mathbb Z$，前两项之和为 $-1$，后两项之和为 $1$，总和仍为 $0$。所以只有 $0$。\end{solution}
\begin{solution}第 4 题：写 $x=[x]+\{x\}$，整数部分可以提出，得所给公式。若 $0\le\{x\}<\frac12$，差为 $0$；若 $\frac12\le\{x\}<1$，差为 $1$。\end{solution}
\begin{solution}第 5 题：设 $x=n+t$。按 $t$ 落在 $[0,\frac14)$、$[\frac14,\frac12)$、$[\frac12,\frac34)$、$[\frac34,1)$ 分段，原式依次为 $4n,4n+1,4n+2,4n+3$，均等于 $[4x]$。\end{solution}
\begin{solution}第 6 题：写 $x=qn+r+t$，其中 $q\in\mathbb Z$，$r\in\{0,1,\ldots,n-1\}$，$0\le t<1$。每项为 $q+[\frac{r+t+k}{n}]$，后面的整数部分恰有 $r$ 个为 $1$，故总和为 $nq+r=[x]$。\end{solution}

\subsection{方程与不等式}
\begin{solution}第 7 题：设 $[x]=n$，则 $x=2n-1$。它是整数，所以 $[x]=x$，即 $n=2n-1$，得 $n=1$、$x=1$。\end{solution}
\begin{solution}第 8 题：设 $x=n+t$。当 $0\le t<\frac12$ 时，左边为 $2n$，不能等于 $5$；当 $\frac12\le t<1$ 时，左边为 $2n+1$，故 $n=2$。解集为 $[\frac52,3)$。\end{solution}
\begin{solution}第 9 题：设 $x=n+t$，则 $2n+[2t]=n+3$，即 $n+[2t]=3$。若 $t<\frac12$，得 $n=3$；若 $t\ge\frac12$，得 $n=2$。故
\[
x\in[3,\tfrac72)\cup[\tfrac52,3)=[\tfrac52,\tfrac72).
\]\end{solution}
\begin{solution}第 10 题：右边必须在 $[0,1)$，所以 $-1\le x<2$。分 $[-1,0)$、$[0,1)$、$[1,2)$ 讨论。代入 $\{x\}=x+1,x,x-1$，分别得到 $x=-1$、$x=\frac12$、$x=2$；最后一个不在区间中。解为 $-1,\frac12$。\end{solution}
\begin{solution}第 11 题：$[3x-2]>1$ 等价于 $3x-2\ge2$，而 $[3x-2]\le6$ 等价于 $3x-2<7$。故 $x\in[\frac43,3)$。\end{solution}
\begin{solution}第 12 题：设 $x=n+t$。条件为 $n<3t$。因 $0\le3t<3$：若 $n\le-1$，全部成立；若 $n=0,1,2$，分别有 $t>0,\frac13,\frac23$；若 $n\ge3$，无解。故
\[
x\in(-\infty,0)\cup(0,1)\cup(\tfrac43,2)\cup(\tfrac83,3).
\]\end{solution}

\subsection{证明与综合提高}
\begin{solution}第 13 题：
\[
[x+y]=[x]+[y]+[\{x\}+\{y\}].
\]
最后一项为 $0$ 或 $1$。左等号在 $\{x\}+\{y\}<1$ 时成立，右等号在 $\{x\}+\{y\}\ge1$ 时成立。\end{solution}
\begin{solution}第 14 题：设 $x=m+t$，$0\le t<1$。右边为
\[
nm+\sum_{k=0}^{n-1}[t+k/n].
\]
设 $r=[nt]$，恰有 $r$ 项满足 $t+k/n\ge1$，所以右边为 $nm+r=[nx]$。\end{solution}
\begin{solution}第 15 题：设 $[x]=qn+r$，其中 $0\le r<n$，再写 $x=[x]+t$，$0\le t<1$。则
\[
\frac{x}{n}=q+\frac{r+t}{n},\qquad0\le\frac{r+t}{n}<1,
\]
两边的整数部分均为 $q$。\end{solution}
\begin{solution}第 16 题：由
\[
x+y+z=[x]+[y]+[z]+\{x\}+\{y\}+\{z\}
\]
知三个小数部分之和是整数；它又属于 $[0,3)$，故只能为 $0,1,2$。\end{solution}
\begin{solution}第 17 题：左边是整数，所以 $7x\in\mathbb Z$。另一方面写 $x=n+t$，等式化为
\[
[2t]+[4t]=7t.
\]
按分点 $0,\frac14,\frac12,\frac34,1$ 检查：四段左边依次为 $0,1,3,4$，相应候选为 $t=0,\frac17,\frac37,\frac47$。后三个候选均不在产生它们的区间中，只有 $t=0$ 成立。故 $x\in\mathbb Z$。\end{solution}
\begin{solution}第 18 题：$n=1$ 显然成立。若 $n\ge2$，取 $x=\frac1n$，则 $[nx]=1$，而 $n[x]=0$，矛盾。因此只有 $n=1$。\end{solution}
\end{document}
